\chapter{Réalisation de la plateforme : sécurité, fiches de paie et virements de masse}
\label{ch:realisation}

Ce chapitre présente la réalisation des trois premiers sprints, qui construisent le socle fonctionnel de SEPAPaie. Le premier met en place les comptes, l'authentification, les sociétés débitrices et le journal d'audit ; le deuxième livre la saisie des fiches de paie, la gestion des schémas XSD et la génération du fichier \texttt{pain.001} ; le troisième assemble ces briques dans un parcours de virement de masse avec contrôles métier et circuit de validation. Chaque sprint est présenté selon la même démarche : conception, puis réalisation, puis écrans obtenus. Les tests automatisés des trois sprints sont regroupés en fin de chapitre.

\chapter{Sprint 1 : socle sécurisé et administration}
\label{ch:sprint1}
\startcontents[chapters]
\chapterplan

\section{Introduction}

Le premier sprint pose les fondations de la plateforme. Avant de manipuler des données de paie, SEPAPaie doit savoir qui l'utilise, ce que chacun a le droit de faire et comment conserver la trace de ses actions. Planifié sur quatre semaines (tableau~\ref{tab:sprints}), il couvre la création et l'approbation des comptes, l'authentification, la gestion des sociétés débitrices, leur affectation aux Experts RH et le journal d'audit.

\section{Backlog du sprint}

Le backlog du sprint reprend les six premiers récits utilisateur du backlog produit (tableau~\ref{tab:backlog}). Le tableau~\ref{tab:backlog-s1} précise les tâches associées à chacun.

\begin{table}[H]
\centering
\footnotesize
\renewcommand{\arraystretch}{1.1}
\tablestyle
\begin{tabularx}{\textwidth}{|>{\centering\arraybackslash}p{0.9cm}|>{\raggedright\arraybackslash}p{5.2cm}|Y|}
\hline
\rowcolor{appMarine}\entete{N\textsuperscript{o}} & \entete{Récit utilisateur} & \entete{Principales tâches} \\
\hline
US1 & Créer un compte avec ma photo pour demander un accès & Formulaire d'inscription avec capture du visage ; contrôle d'unicité de l'identifiant et de l'adresse électronique ; calcul et conservation de l'empreinte du visage \\
\hline
US2 & M'authentifier par mot de passe et code à usage unique, ou par reconnaissance faciale & Vérification du mot de passe ; envoi, contrôle et expiration du code ; comparaison de visage ; émission du jeton de session \\
\hline
US3 & Valider les inscriptions et gérer les comptes et les rôles & Liste des demandes en attente ; approbation et rejet ; création directe d'un compte ; désactivation et réactivation \\
\hline
US4 & Gérer les sociétés débitrices et leurs coordonnées bancaires & Création et modification d'une société ; contrôle de l'IBAN et du BIC ; masquage de l'IBAN à l'affichage \\
\hline
US5 & Affecter les sociétés aux Experts RH & Affectation et retrait d'une société ; conservation de l'historique des affectations \\
\hline
US6 & Consulter le journal d'audit & Enregistrement des événements ; consultation réservée à l'administrateur, du plus récent au plus ancien \\
\hline
\end{tabularx}
\caption{Backlog du sprint 1}
\label{tab:backlog-s1}
\end{table}

\section{Conception}

\subsection{Diagramme de cas d'utilisation}

La figure~\ref{fig:uc-s1} détaille les cas d'utilisation de ce sprint. Deux acteurs système y apparaissent : le service Python, qui calcule et compare les empreintes de visage, et la messagerie, qui achemine le code à usage unique. L'administrateur gère les comptes, les sociétés et consulte le journal d'audit ; le Manager affecte les sociétés aux Experts RH ; tout utilisateur peut s'authentifier et gérer son propre visage de référence.

\begin{figure}[H]
\centering
\includegraphics[height=0.38\textheight,width=\textwidth,keepaspectratio]{fig4_1_cas_utilisation_sprint1}
\caption{Cas d'utilisation du sprint 1}
\label{fig:uc-s1}
\end{figure}

\subsection{Inscription et approbation des comptes}

Un visiteur ne peut pas utiliser la plateforme dès son inscription. Il renseigne son identité, choisit le profil demandé et capture une photo de son visage. Le backend refuse la demande si l'identifiant ou l'adresse électronique existent déjà, ou si le profil n'est pas attribuable. Il sollicite ensuite le service Python pour obtenir l'empreinte numérique du visage. Seule cette empreinte est conservée en base ; la photo elle-même n'est pas enregistrée. Le compte est alors créé \textbf{inactif} et \textbf{non approuvé}, avec un mot de passe haché. La figure~\ref{fig:seq-inscription} illustre ce parcours.

\begin{figure}[H]
\centering
\includegraphics[height=0.55\textheight,width=\textwidth,keepaspectratio]{fig4_3_sequence_inscription}
\caption{Séquence d'inscription et d'approbation d'un compte}
\label{fig:seq-inscription}
\end{figure}

L'administrateur consulte les demandes en attente et les approuve ou les rejette telles qu'elles ont été formulées. L'approbation active le compte ; le rejet supprime la demande. L'administrateur peut aussi créer directement un compte, désactiver un compte existant ou le réactiver. La figure~\ref{fig:etats-compte} résume le cycle de vie d'un compte.

\begin{figure}[H]
\centering
\includegraphics[width=\textwidth]{fig4_4_etats_compte}
\caption{Cycle de vie d'un compte utilisateur}
\label{fig:etats-compte}
\end{figure}

\subsection{Authentification adaptative}

SEPAPaie propose deux méthodes de connexion. Avec la première, l'utilisateur saisit son identifiant et son mot de passe, puis un code à usage unique envoyé par courriel. Avec la seconde, il présente son visage devant la caméra. La plateforme est dite \textbf{adaptative} parce que le backend, et non l'interface, décide si un second facteur est exigé : un mot de passe correct ne suffit jamais à ouvrir une session, alors qu'un visage reconnu constitue lui-même le facteur d'authentification. Cette règle est centralisée en un point unique du code, ce qui évite qu'une modification de l'interface ne la contourne.

La figure~\ref{fig:seq-auth} décrit la connexion par mot de passe. Le code est généré aléatoirement, conservé sous forme hachée et soumis à des limites décrites dans le tableau~\ref{tab:otp}.

\begin{figure}[H]
\centering
\includegraphics[height=0.58\textheight,width=\textwidth,keepaspectratio]{fig4_2_sequence_authentification}
\caption{Séquence d'authentification par mot de passe et code à usage unique}
\label{fig:seq-auth}
\end{figure}

\begin{table}[H]
\centering
\small
\renewcommand{\arraystretch}{1.25}
\tablestyle
\begin{tabularx}{\textwidth}{|>{\raggedright\arraybackslash}p{5cm}|Y|}
\hline
\rowcolor{appMarine}\entete{Paramètre} & \entete{Valeur par défaut} \\
\hline
Format du code & Six chiffres tirés par un générateur aléatoire cryptographique \\
\hline
Conservation du code & Sous forme hachée, jamais en clair \\
\hline
Durée de validité & Cinq minutes \\
\hline
Tentatives de saisie & Cinq au maximum ; au-delà, une nouvelle connexion est exigée \\
\hline
Renvois du code & Trois au maximum, avec un délai de trente secondes entre deux envois \\
\hline
\end{tabularx}
\caption{Paramètres du code à usage unique}
\label{tab:otp}
\end{table}

Ces valeurs sont configurables par variables d'environnement. Le principe retenu est de ne jamais accorder de session sur la seule foi d'un mot de passe : même si le réglage général du code à usage unique est désactivé, le parcours par mot de passe continue de l'exiger.

Pour la reconnaissance faciale, le backend transmet l'image capturée et l'empreinte de référence au service Python, qui détecte le visage, calcule son empreinte et mesure sa similarité avec la référence. Si elle dépasse le seuil recommandé par la bibliothèque, le visage est reconnu et le jeton est émis sans code. Dans tous les autres cas (visage refusé, service indisponible, erreur inattendue), aucun jeton n'est délivré et l'interface invite l'utilisateur à se connecter par mot de passe. L'indisponibilité du service Python ne peut donc jamais ouvrir un accès.

\subsection{Sécurité de la session et contrôle des droits}

Après une authentification réussie, le backend émet un jeton \textbf{JWT} signé, valable vingt-quatre heures. À chaque requête, un filtre en extrait l'identité, puis recharge l'utilisateur depuis la base : le rôle n'est donc pas lu dans le jeton, ce qui garantit que les droits appliqués sont ceux enregistrés au moment de l'appel. Aucune session n'est conservée côté serveur. Les mots de passe sont hachés avec l'algorithme BCrypt.

Les droits sont ensuite appliqués côté serveur. Les routes d'authentification sont publiques ; toutes les autres exigent un utilisateur authentifié. Chaque point d'accès est de plus protégé par une annotation qui reprend les expressions de droits centralisées dans une classe unique, conformément au tableau~\ref{tab:droits}. L'interface masque les menus non autorisés, mais cette présentation ne remplace pas le contrôle du serveur.

Le modèle de rôles a évolué en cours de projet : un premier découpage en quatre profils a été remplacé par les trois profils actuels (administrateur, Manager, Expert RH). Une migration exécutée au démarrage convertit les comptes existants vers les nouveaux rôles de façon idempotente et trace l'opération dans le journal d'audit. Les profils obsolètes ne sont plus attribuables.

\subsection{Gestion des utilisateurs}

La gestion des comptes est réservée à l'administrateur. Il y retrouve les demandes en attente, signalées par un message, ainsi que les comptes actifs et désactivés. Désactiver un compte empêche toute nouvelle connexion sans supprimer ses données, ce qui préserve l'historique des actions de la personne. Un utilisateur peut par ailleurs enregistrer, remplacer ou supprimer son visage de référence depuis son profil ; sans visage enregistré, seule la connexion par mot de passe lui est proposée.

\subsection{Sociétés débitrices}

Une société débitrice est caractérisée par un code, une raison sociale, un IBAN et un BIC. Lors de la création ou de la modification, l'IBAN est vérifié par le contrôle modulo~97 et le BIC par son format. Un enregistrement invalide est refusé avec un message explicite. Le code d'une société n'est pas modifiable, car il relie la société aux fiches de paie et aux fichiers générés. Une société n'est pas supprimée mais désactivée.

L'IBAN est masqué dans les listes : seuls les quatre premiers et les quatre derniers caractères apparaissent. Lors d'une modification, un champ IBAN laissé vide conserve la valeur existante, ce qui évite de réafficher la donnée complète. L'IBAN complet n'est communiqué qu'aux profils qui préparent les virements, puisque la génération du fichier en a besoin. Les coordonnées d'une société inconnue, inactive ou incomplète sont signalées comme telles afin d'expliquer pourquoi la préparation d'un fichier est impossible.

\subsection{Affectation des sociétés aux Experts RH}

Le Manager affecte chaque société à un Expert RH actif. Une société n'a qu'un expert à la fois, alors qu'un expert peut gérer plusieurs sociétés. Lorsqu'une société est réaffectée, l'affectation précédente est clôturée avec sa date de fin au lieu d'être écrasée : l'historique reste consultable et chaque changement est tracé dans le journal d'audit. Cette affectation définit le périmètre de travail de l'Expert RH, sur lequel s'appuient les contrôles d'accès par société présentés dans la suite du rapport.

\subsection{Journal d'audit}

Le journal enregistre, pour chaque événement, l'horodatage, l'utilisateur, le type d'action, un détail, l'adresse IP et le résultat. Dès ce sprint, il retient les étapes de connexion (code envoyé, connexion réussie, échec), les migrations de rôles et les changements d'affectation. Il est enrichi aux sprints suivants par les opérations sur les fiches et les fichiers. Seul l'administrateur peut le consulter. L'interface affiche les cent événements les plus récents, du plus récent au plus ancien, à partir d'une interface de programmation qui accepte la pagination, et aucune fonction ne permet d'en modifier le contenu. L'enregistrement des événements reste en revanche commun à tous les profils.

\subsection{Interfaces réalisées}

Les écrans suivants proviennent de l'application exécutée localement avec des comptes de démonstration. La figure~\ref{fig:ecrans-connexion} présente la connexion, dans laquelle le choix entre mot de passe et reconnaissance faciale est proposé sur le même écran, puis l'étape de saisie du code reçu. L'adresse de destination y est masquée.

\begin{figure}[H]
\centering
\begin{subfigure}[t]{0.36\textwidth}
\centering
\includegraphics[width=\textwidth]{fig4_login}
\caption{Saisie des identifiants}
\end{subfigure}
\hspace{0.05\textwidth}
\begin{subfigure}[t]{0.36\textwidth}
\centering
\includegraphics[width=\textwidth]{fig4_otp}
\caption{Saisie du code reçu}
\end{subfigure}
\caption{Écrans de connexion}
\label{fig:ecrans-connexion}
\end{figure}

La figure~\ref{fig:ecran-societes} présente la liste des sociétés débitrices, dont l'IBAN est masqué, et la figure~\ref{fig:ecran-audit} le journal d'audit, où chaque connexion laisse la trace de l'envoi du code puis de sa vérification.

\begin{figure}[H]
\centering
\includegraphics[width=0.9\textwidth]{fig4_societes}
\caption{Sociétés débitrices (profil administrateur)}
\label{fig:ecran-societes}
\end{figure}

\begin{figure}[H]
\centering
\includegraphics[width=0.9\textwidth]{fig4_audit}
\caption{Journal d'audit (profil administrateur)}
\label{fig:ecran-audit}
\end{figure}

\section{Conclusion}

Le sprint 1 a livré le socle des suivants : comptes soumis à approbation, authentification qui n'accorde jamais d'accès sur la seule foi d'un mot de passe, droits appliqués côté serveur, sociétés débitrices contrôlées, affectations historisées et journal d'audit. Le chapitre suivant présente le sprint 2, consacré aux données de paie, aux schémas XSD et à la génération du fichier SEPA.


\chapter{Sprint 2 : données de paie, schémas XSD et génération du fichier SEPA}
\label{ch:sprint2}
\startcontents[chapters]
\chapterplan

\section{Introduction}

Le deuxième sprint construit le c\oe ur fonctionnel de la plateforme : transformer des fiches de paie en un fichier \texttt{pain.001} conforme à un schéma choisi. Il couvre la saisie et l'approbation des fiches, la gestion des schémas XSD, la génération du message et sa validation, avec des messages d'erreur compréhensibles par un utilisateur non spécialiste. Planifié sur quatre semaines (tableau~\ref{tab:sprints}), il livre le noyau technique sur lequel s'appuie, au sprint suivant, le parcours complet de virement de masse.

\section{Backlog du sprint}

Le sprint reprend les récits US7 à US11 du backlog produit (tableau~\ref{tab:backlog}). Le tableau~\ref{tab:backlog-s2} précise les tâches associées.

\begin{table}[H]
\centering
\footnotesize
\renewcommand{\arraystretch}{1.1}
\tablestyle
\begin{tabularx}{\textwidth}{|>{\centering\arraybackslash}p{0.9cm}|>{\raggedright\arraybackslash}p{5.2cm}|Y|}
\hline
\rowcolor{appMarine}\entete{N\textsuperscript{o}} & \entete{Récit utilisateur} & \entete{Principales tâches} \\
\hline
US7 & Saisir des fiches de paie pour mes sociétés et les supprimer & Formulaire de saisie avec choix de la société ; contrôle d'accès par société ; suppression avec confirmation \\
\hline
US8 & Approuver les fiches de paie & Approbation par le Manager ; vérification des données avant approbation ; traçabilité \\
\hline
US9 & Importer et gérer les schémas XSD & Import, mise à jour et suppression de schémas ; droits réservés à l'administrateur et au Manager \\
\hline
US10 & Générer un fichier \texttt{pain.001} selon le schéma choisi & Lecture de la version et de la structure du schéma ; construction du message ; adaptation aux versions \\
\hline
US11 & Valider le fichier contre le XSD et en prévisualiser le contenu & Validation avec collecte de toutes les erreurs ; messages en clair ; aperçu du XML ; enregistrement dans l'historique \\
\hline
\end{tabularx}
\caption{Backlog du sprint 2}
\label{tab:backlog-s2}
\end{table}

\section{Conception}

\subsection{Diagramme de cas d'utilisation}

La figure~\ref{fig:uc-s2} présente les cas d'utilisation du sprint. L'Expert RH saisit et supprime les fiches, puis génère, prévisualise et télécharge le fichier. Le Manager approuve les fiches, tandis que la gestion des schémas XSD est partagée entre le Manager et l'administrateur. La génération du fichier inclut deux étapes systématiques : l'analyse du schéma choisi et la validation du résultat contre ce schéma.

\begin{figure}[H]
\centering
\includegraphics[width=\textwidth,height=0.36\textheight,keepaspectratio]{fig5_1_cas_utilisation_sprint2}
\caption{Cas d'utilisation du sprint 2}
\label{fig:uc-s2}
\end{figure}

\subsection{Cycle de vie d'une fiche de paie}

Une fiche de paie correspond à un virement individuel destiné à un salarié. Elle est saisie à l'état \textbf{brouillon}, puis approuvée par le Manager. Seules les fiches approuvées peuvent entrer dans un fichier SEPA. Cette séparation entre la personne qui saisit et celle qui approuve répète, à l'échelle de la fiche, le principe des quatre yeux présenté au chapitre~\ref{ch:besoins}. La figure~\ref{fig:etats-fiche} résume ce cycle de vie.

\begin{figure}[H]
\centering
\includegraphics[width=\textwidth]{fig5_2_etats_fiche}
\caption{Cycle de vie d'une fiche de paie}
\label{fig:etats-fiche}
\end{figure}

\subsection{Chaîne de génération et de validation}

La génération du fichier suit une chaîne en plusieurs étapes, illustrée par la figure~\ref{fig:generation}. Le schéma choisi n'est pas seulement un outil de contrôle appliqué à la fin : il guide aussi la construction du message, dont certains éléments dépendent de la version. La plateforme analyse donc d'abord le schéma, construit le document en conséquence, puis le valide et présente le résultat. En cas de non-conformité, le fichier ne peut être ni téléchargé ni soumis.

\begin{figure}[H]
\centering
\begin{tikzpicture}[
  etape/.style={draw=appMarine, fill=appClair, rounded corners=3pt, align=center, text width=4.4cm, minimum height=1.25cm, font=\small, execute at begin node={\hyphenpenalty=10000\relax}},
  refus/.style={draw=acteurSysteme, fill=orange!8, dashed, rounded corners=3pt, align=center, text width=4.4cm, minimum height=1.0cm, font=\small, execute at begin node={\hyphenpenalty=10000\relax}},
  choix/.style={draw=acteurSysteme, fill=orange!8, diamond, aspect=2.6, inner sep=1pt, align=center, font=\small},
  fleche/.style={-{Stealth[length=2.4mm]}, appMarine, thick}]
\node[etape] (a) {1. Sélection des fiches\\approuvées et choix du schéma};
\node[etape, right=0.7cm of a] (b) {2. Analyse du schéma\\(version, structure)};
\node[etape, right=0.7cm of b] (c) {3. Construction du\\document XML};
\node[refus, above=0.9cm of b] (x) {Refus avec\\message explicatif};
\node[etape, below=0.9cm of c] (d) {4. Validation contre le XSD\\(toutes les erreurs)};
\node[etape, below=0.9cm of b] (e) {5. Messages d'erreur\\en clair (ligne, balise)};
\node[choix, below=0.9cm of a] (f) {Fichier conforme ?};
\node[etape, below=1.0cm of f] (g) {Enregistrement (statut\\Brouillon), contenu figé ;\\téléchargement possible};
\node[refus, right=0.7cm of g] (h) {Erreurs présentées ;\\téléchargement et soumission\\refusés};
\draw[fleche] (a) -- (b);
\draw[fleche] (b) -- (c);
\draw[fleche, acteurSysteme] (b) -- (x) node[midway, right, font=\scriptsize, text=black] {schéma incompatible};
\draw[fleche] (c) -- (d);
\draw[fleche] (d) -- (e);
\draw[fleche] (e) -- (f);
\draw[fleche] (f.south) -- (g.north) node[midway, right, font=\scriptsize] {oui};
\draw[fleche] (f.south) -- ++(0,-0.5) -| (h.north) node[pos=0.2, above, font=\scriptsize] {non};
\end{tikzpicture}
\caption{Chaîne de génération et de validation d'un fichier \texttt{pain.001}}
\label{fig:generation}
\end{figure}

\subsection{Fiches de paie et approbation}

Une fiche de paie comporte l'identifiant du salarié, le nom du bénéficiaire, son IBAN, un BIC facultatif, le montant, la devise (l'euro par défaut), un libellé et le code de la société débitrice. L'Expert RH la saisit pour l'une des sociétés qui lui sont affectées : le serveur refuse la saisie si la société est inconnue, inactive ou non affectée à l'utilisateur, et la liste proposée dans le formulaire est déjà limitée à ses sociétés. Le montant doit être strictement positif. Dans les listes, l'IBAN est masqué et la création comme la suppression d'une fiche sont tracées dans le journal d'audit.

L'approbation est réservée au Manager. Avant de la prononcer, le serveur vérifie les données de la fiche : IBAN valide (format et clé de contrôle), BIC bien formé s'il est renseigné et montant positif. Si l'une de ces vérifications échoue, la fiche reste en brouillon et le message indique ce qu'il faut corriger. Ces règles, communes à plusieurs étapes de la plateforme, sont détaillées au chapitre~\ref{ch:sprint3}. Dans le cadre du projet, les fiches proviennent de cette saisie unitaire ou de jeux de données synthétiques chargés en base pour la démonstration et l'évaluation ; un import automatique depuis le SIRH n'est pas réalisé.

\subsection{Gestion des schémas XSD}

Un schéma XSD est importé sous forme de fichier. Il est conservé avec son nom, son contenu, le nom de l'utilisateur qui l'a importé et sa date d'import. Importer à nouveau un fichier de même nom crée une nouvelle version, et la mise à jour d'un schéma existant en incrémente également le numéro de version. La suppression est possible. L'import, la mise à jour et la suppression sont réservés à l'administrateur et au Manager, et chacune de ces opérations est enregistrée dans le journal d'audit. Tous les utilisateurs authentifiés peuvent en revanche consulter la liste des schémas disponibles.

L'import n'impose pas qu'un fichier soit un schéma \texttt{pain.001}. La compatibilité est vérifiée au moment de l'utilisation : un schéma qui ne décrit pas l'élément \texttt{CstmrCdtTrfInitn} est signalé comme incompatible et ne peut pas servir à générer un fichier. Ce choix permet de conserver d'autres schémas, par exemple d'autres versions, sans bloquer leur import.

\subsection{Analyse dynamique du schéma}

Pour construire un message cohérent avec le schéma choisi, la plateforme en lit la structure. Elle extrait d'abord la version et l'espace de noms, puis parcourt l'arbre des éléments à partir de \texttt{CstmrCdtTrfInitn}. Pour chaque élément, elle relève s'il est obligatoire, s'il peut se répéter, son type et ses contraintes : longueur maximale ou minimale, format, nombre de chiffres, liste de valeurs autorisées. Le test sur le schéma officiel \texttt{pain.001.001.03} confirme par exemple que l'identifiant du message est limité à 35 caractères et le nom d'une partie à 140.

Cette lecture est protégée de deux manières. Elle est bornée en profondeur et en nombre d'éléments, et elle détecte les définitions circulaires, afin qu'un schéma volumineux ou mal conçu ne bloque pas le serveur. Elle interdit aussi les déclarations de type de document et l'accès aux ressources externes, ce qui écarte les attaques par entités externes déjà évoquées pour les besoins non fonctionnels.

\subsection{Génération du fichier \texttt{pain.001}}

Le document est construit avec l'interface standard de manipulation d'arbres XML, puis sérialisé avec indentation. Il comporte un en-tête de groupe, un lot de paiement unique pour la société débitrice, puis une instruction par fiche. Le nombre de virements et la somme de contrôle sont calculés à partir des fiches sélectionnées. Le lot porte le niveau de service \texttt{SEPA} et la catégorie \texttt{SALA}, qui qualifie un virement de salaire. Les coordonnées de la société débitrice proviennent de sa fiche, présentée au chapitre~\ref{ch:sprint1}.

Plusieurs versions du message coexistent. Pour les versions \texttt{pain.001.001.09} à \texttt{.11}, le code de la banque est porté par l'élément \texttt{BICFI} et la date d'exécution est encapsulée dans un élément \texttt{Dt} ; pour la version \texttt{.03}, ils utilisent respectivement \texttt{BIC} et une date simple. La version est lue dans le schéma choisi et le générateur adapte le message en conséquence.

Les valeurs saisies ne respectent pas toujours les contraintes du standard. Le générateur les prépare donc selon les règles du tableau~\ref{tab:preparation-valeurs}. Lorsqu'un champ facultatif est vide, comme le nom d'une partie ou le BIC, et que le schéma l'autorise, la balise est omise et un avertissement non bloquant est produit. Un champ obligatoire vide reste en revanche une erreur de validation.

\begin{table}[H]
\centering
\small
\renewcommand{\arraystretch}{1.2}
\tablestyle
\begin{tabularx}{\textwidth}{|>{\raggedright\arraybackslash}p{4.2cm}|Y|}
\hline
\rowcolor{appMarine}\entete{Donnée} & \entete{Préparation appliquée} \\
\hline
Noms et libellés & Suppression des accents, remplacement des caractères hors du jeu autorisé, puis troncature : 70 caractères pour un nom, 140 pour le libellé \\
\hline
Identifiants du message et des virements & Identifiant unique composé de la date et d'une partie aléatoire ; référence de bout en bout construite à partir du matricule et de la période ; troncature à 35 caractères \\
\hline
Montants & Deux décimales, avec le point comme séparateur \\
\hline
IBAN et BIC & Suppression des espaces et passage en majuscules \\
\hline
Date d'exécution & Valeur demandée, ou à défaut le jour courant augmenté de deux jours \\
\hline
Libellé du virement & \og Salaire\fg{} suivi du mois, de l'année et du matricule \\
\hline
\end{tabularx}
\caption{Préparation des valeurs lors de la génération}
\label{tab:preparation-valeurs}
\end{table}

\subsection{Validation XSD et messages d'erreur}

Le document généré est validé contre le schéma choisi. Par défaut, un validateur s'arrête à la première erreur et produit des messages techniques en anglais, difficiles à exploiter. SEPAPaie collecte donc \textbf{toutes} les erreurs en une seule passe. Pour chacune, elle conserve la ligne, la colonne et la balise concernée, puis la traduit en un message en français qui indique la correction à apporter. Plusieurs diagnostics situés au même endroit sont regroupés, et un message particulier distingue l'IBAN du débiteur de celui d'un bénéficiaire.

Le tableau~\ref{tab:erreurs-xsd} donne des exemples de messages obtenus. L'utilisateur n'a ainsi pas besoin de connaître le vocabulaire du standard pour comprendre ce qu'il doit corriger.

\begin{table}[H]
\centering
\small
\renewcommand{\arraystretch}{1.2}
\tablestyle
\begin{tabularx}{\textwidth}{|>{\raggedright\arraybackslash}p{3.6cm}|Y|}
\hline
\rowcolor{appMarine}\entete{Situation} & \entete{Message restitué à l'utilisateur (reformulé)} \\
\hline
Valeur trop longue & Le nom est trop long ; raccourcissez-le au nombre de caractères maximal du schéma \\
\hline
IBAN au format incorrect & L'IBAN du compte concerné est incorrect ; il doit comporter deux lettres de pays, deux chiffres de contrôle, puis le numéro de compte, sans espace \\
\hline
Date invalide & La date d'exécution est invalide ; utilisez une date réelle au format année-mois-jour \\
\hline
Montant invalide & Le montant est invalide ; utilisez un nombre avec un point décimal, sans devise ni séparateur de milliers \\
\hline
Balise manquante ou mal placée & Une balise est manquante, inattendue ou placée dans le mauvais ordre ; le message précise la balise attendue \\
\hline
\end{tabularx}
\caption{Exemples de messages d'erreur de validation XSD}
\label{tab:erreurs-xsd}
\end{table}

\subsection{Aperçu, enregistrement et téléchargement}

L'Expert RH peut d'abord prévisualiser le fichier sans l'enregistrer. L'aperçu présente le XML généré, le résultat de la validation et les erreurs éventuelles. Lorsque le fichier est conforme, il peut être enregistré dans l'historique. L'enregistrement conserve le contenu XML, le schéma utilisé et sa version, la société et la période, les matricules et les identifiants des fiches incluses, le nombre de virements, la somme de contrôle, les coordonnées du débiteur, le résultat de la validation, l'auteur et la date. Le fichier reçoit alors le statut \og Brouillon\fg{}.

Le contenu XML n'est jamais modifié après l'enregistrement : un téléchargement ultérieur restitue exactement le même document, ce qui permet de le rapprocher du fichier remis à la banque. Un fichier non conforme ne peut pas être téléchargé. Le téléchargement d'un fichier qui n'a pas encore été validé par le Manager donne un nom préfixé par \og NON\_VALIDE\fg{}, afin d'éviter qu'il soit confondu avec un fichier validé. Chaque téléchargement est tracé. L'historique peut être filtré par société, période, conformité et statut, et un Expert RH n'y voit que les fichiers de ses sociétés. Le parcours complet, qui enchaîne contrôles métier, analyse par l'IA et validation XSD en conservant aussi les résultats non conformes pour consultation, est présenté au chapitre~\ref{ch:sprint3}.

\subsection{Interfaces réalisées}

La figure~\ref{fig:ecran-saisie} montre la saisie d'une fiche par un Expert RH : la société est choisie parmi celles qui lui sont affectées. Le tableau indique la société sous le matricule, le statut de la fiche et, dans les colonnes de droite, les informations du module d'IA présenté au chapitre~\ref{ch:sprint4}.

\begin{figure}[H]
\centering
\includegraphics[width=0.95\textwidth]{fig5_saisie_fiche}
\caption{Saisie d'une fiche de paie (profil Expert RH)}
\label{fig:ecran-saisie}
\end{figure}

Les figures~\ref{fig:ecran-xml} et~\ref{fig:ecran-validation} montrent le résultat d'un fichier généré. L'onglet XML présente le document avec une coloration syntaxique ; les lignes en erreur ou en avertissement y sont mises en évidence. L'onglet de validation indique la conformité au schéma et liste, le cas échéant, les erreurs avec leur ligne et leur balise.

\begin{figure}[H]
\centering
\includegraphics[width=0.95\textwidth]{fig5_xml}
\caption{Aperçu du fichier XML généré}
\label{fig:ecran-xml}
\end{figure}

\begin{figure}[H]
\centering
\includegraphics[width=0.95\textwidth]{fig5_validation_xsd}
\caption{Résultat de la validation XSD}
\label{fig:ecran-validation}
\end{figure}

\section{Conclusion}

Le sprint 2 a livré le noyau technique de la plateforme : fiches de paie rattachées à une société et approuvées avant tout usage, schémas XSD gérés par les profils autorisés, fichier \texttt{pain.001} construit selon la version du schéma, et validation qui restitue toutes les erreurs en langage clair. Le chapitre suivant présente le sprint 3, qui enchaîne ces briques dans un parcours de virement de masse avec contrôles métier et circuit de validation.


\chapter{Sprint 3 : virements de masse, contrôles métier et circuit de validation}
\label{ch:sprint3}
\startcontents[chapters]
\chapterplan

\section{Introduction}

Les deux premiers sprints ont fourni les comptes, les sociétés, les fiches de paie et la génération d'un fichier conforme. Le troisième sprint assemble ces briques dans un parcours complet, de la sélection des paiements à la décision du Manager. Il ajoute ce qui distingue SEPAPaie d'un simple générateur de fichiers : des contrôles métier explicites, un suivi des anomalies assorti de recommandations et un circuit de validation fondé sur le principe des quatre yeux. Ce sprint, planifié sur quatre semaines (tableau~\ref{tab:sprints}), concerne les récits US12 à US16.

\section{Backlog du sprint}

\begin{table}[H]
\centering
\footnotesize
\renewcommand{\arraystretch}{1.1}
\tablestyle
\begin{tabularx}{\textwidth}{|>{\centering\arraybackslash}p{0.9cm}|>{\raggedright\arraybackslash}p{5.2cm}|Y|}
\hline
\rowcolor{appMarine}\entete{N\textsuperscript{o}} & \entete{Récit utilisateur} & \entete{Principales tâches} \\
\hline
US12 & Contrôler un virement de masse par des règles métier & Parcours de préparation ; règles sur l'IBAN, le BIC, le montant et les doublons ; restitution des erreurs par salarié \\
\hline
US13 & Consulter et traiter les anomalies avec des recommandations & Registre des anomalies ; statuts et commentaire ; catalogue de recommandations \\
\hline
US14 & Soumettre un fichier au Manager & Conditions de soumission ; statut du fichier ; traçabilité \\
\hline
US15 & Valider ou rejeter un fichier soumis, avec un motif & Liste des fichiers à valider ; principe des quatre yeux ; motif obligatoire au rejet \\
\hline
US16 & Consulter l'historique des fichiers & Historique filtrable ; limitation aux sociétés de l'Expert RH ; vue de supervision de l'administrateur \\
\hline
\end{tabularx}
\caption{Backlog du sprint 3}
\label{tab:backlog-s3}
\end{table}

\section{Conception}

\subsection{Diagramme de cas d'utilisation}

La figure~\ref{fig:uc-s3} présente les cas d'utilisation du sprint. L'Expert RH prépare et contrôle le virement, puis soumet le fichier. Le Manager le valide ou le rejette. Les deux profils métier consultent le résultat, l'historique et traitent les anomalies, alors que l'administrateur se limite à une supervision en lecture seule. L'analyse des montants par l'IA \textbf{étend} le contrôle plutôt que de l'inclure : elle enrichit le résultat lorsqu'elle est disponible, mais son absence n'empêche pas le contrôle d'aboutir.

\begin{figure}[H]
\centering
\includegraphics[height=0.4\textheight,width=\textwidth,keepaspectratio]{fig6_1_cas_utilisation_sprint3}
\caption{Cas d'utilisation du sprint 3}
\label{fig:uc-s3}
\end{figure}

\subsection{Parcours en trois étapes}

Le virement de masse est guidé par un parcours en trois étapes affiché en haut de chaque page : la \textbf{préparation}, le \textbf{contrôle} et la \textbf{validation et le téléchargement}. Cette organisation répond à un besoin de lisibilité : l'utilisateur doit savoir à tout moment quelle vérification a été faite et ce qu'il lui reste à traiter. Chaque étape a donc une page dédiée, et le résultat du contrôle est lui-même présenté en cinq onglets : synthèse, XML, validation XSD, contrôle IA et contrôles métier.

\subsection{Cycle de vie d'un fichier}

Un fichier suit un cycle de vie à cinq statuts, illustré par la figure~\ref{fig:etats-fichier}. Le contrôle produit un fichier \og Brouillon\fg{} s'il est conforme au schéma, ou \og Non conforme\fg{} dans le cas contraire. Seul un brouillon peut être soumis. Le Manager décide ensuite : le fichier devient \og Validé par le Manager\fg{} ou \og Rejeté\fg{}. Un fichier rejeté ou non conforme n'est pas corrigé sur place : l'Expert RH relance un nouveau contrôle depuis la préparation, ce qui crée un nouveau fichier et conserve l'ancien dans l'historique.

\begin{figure}[H]
\centering
\includegraphics[width=\textwidth]{fig6_2_etats_fichier}
\caption{Cycle de vie d'un fichier SEPA}
\label{fig:etats-fichier}
\end{figure}

Il convient de ne pas confondre ce statut avec celui d'une fiche de paie, décrit au chapitre~\ref{ch:sprint2}. L'approbation concerne un paiement individuel ; la validation concerne le lot soumis. La conformité au schéma XSD, enfin, est une propriété technique du fichier distincte de la décision du Manager.

\subsection{Principe des quatre yeux}

La figure~\ref{fig:seq-validation} décrit la soumission puis la décision. Le principe des quatre yeux s'applique à la décision : la personne qui a généré ou soumis un fichier ne peut ni le valider ni le rejeter. Cette règle est contrôlée par le serveur, et l'interface la reflète en désactivant les boutons et en expliquant pourquoi. Elle s'ajoute à la séparation des rôles : seul un Manager peut décider, et seul un Expert RH peut soumettre.

\begin{figure}[H]
\centering
\includegraphics[height=0.5\textheight,width=\textwidth,keepaspectratio]{fig6_3_sequence_validation}
\caption{Séquence de soumission et de décision}
\label{fig:seq-validation}
\end{figure}

\subsection{Préparation d'un virement}

La première page guide l'Expert RH en trois blocs. Il choisit d'abord la \textbf{société débitrice} parmi celles qui lui sont affectées ; l'interface affiche alors sa raison sociale, son BIC et son IBAN masqué, que l'utilisateur peut révéler ponctuellement. Si les coordonnées bancaires de la société sont incomplètes ou si la société est inactive, la génération est bloquée avec un message qui indique ce qu'il faut compléter. Un fichier ne concerne qu'une seule société.

Il renseigne ensuite la période de paie, le schéma XSD et la date d'exécution, proposée par défaut à deux jours ouvrés. Une analyse du schéma choisi peut être consultée. Il examine enfin la liste des paiements approuvés de la société pour la période choisie. Pour chaque paiement, un indicateur affiche le résultat des contrôles de données et, le cas échéant, un signalement de l'IA. Les paiements en erreur sont décochés par défaut, mais restent consultables et peuvent être cochés. L'utilisateur peut filtrer par nom ou matricule et suit en direct le nombre de paiements retenus et leur total (figure~\ref{fig:ecran-preparation}).

Le bouton de lancement reste inactif tant qu'une condition manque, et la page indique laquelle. Une confirmation rappelle le nombre de paiements et le total avant le lancement. Si l'utilisateur quitte la page avec une sélection en cours, l'interface demande confirmation afin d'éviter une perte accidentelle.

\subsection{Lancement du contrôle}

Le lancement déclenche un unique appel au serveur, qui enchaîne dans l'ordre : les contrôles métier, l'analyse par l'IA, la génération du XML puis sa validation contre le schéma. Pendant le traitement, l'interface affiche l'étape en cours. Le résultat est ensuite enregistré dans l'historique, qu'il soit conforme ou non : un fichier non conforme est conservé pour que ses erreurs restent consultables, mais il ne peut être ni téléchargé ni soumis. L'opération est tracée dans le journal d'audit.

Si le service d'IA est désactivé ou indisponible, le traitement se poursuit et le résultat l'indique explicitement (\og Analyse IA indisponible\fg{}). L'absence d'analyse n'est donc jamais confondue avec une absence de risque détecté.

\subsection{Contrôles métier}

Les contrôles métier traduisent des règles explicites sur chaque paiement. Le tableau~\ref{tab:regles-metier} les présente avec leur gravité.

\begin{table}[H]
\centering
\small
\renewcommand{\arraystretch}{1.2}
\tablestyle
\begin{tabularx}{\textwidth}{|>{\raggedright\arraybackslash}p{3.4cm}|Y|>{\centering\arraybackslash}p{2.3cm}|}
\hline
\rowcolor{appMarine}\entete{Contrôle} & \entete{Règle appliquée} & \entete{Gravité} \\
\hline
IBAN & Présence, format général, puis clé de contrôle modulo~97 & Erreur \\
\hline
BIC & Si renseigné : huit ou onze caractères conformes au format attendu & Erreur \\
\hline
Montant & Renseigné et strictement positif & Erreur \\
\hline
Doublon potentiel & Même IBAN et même montant apparaissant plusieurs fois dans la sélection & Avertissement \\
\hline
\end{tabularx}
\caption{Règles de contrôle métier appliquées avant génération}
\label{tab:regles-metier}
\end{table}

Chaque message identifie le paiement par son matricule et indique l'action à mener. Les erreurs sont aussi positionnées sur la ligne du XML correspondante, ce qui permet de passer d'un message à l'élément concerné. Une erreur métier empêche la \textbf{soumission} du fichier, même si le XML est conforme au schéma : la conformité structurelle et l'exactitude des données restent deux vérifications distinctes. L'écran de la figure~\ref{fig:ecran-controles} en donne un exemple.

L'avertissement de doublon ne bloque pas : deux paiements identiques peuvent être légitimes, mais ils méritent une vérification. De même, un montant signalé par l'IA produit un avertissement et non une erreur, comme l'établit le chapitre~\ref{ch:sprint4}.

\subsection{Anomalies et recommandations}

En complément des contrôles lancés sur un fichier, un registre d'anomalies conserve les problèmes détectés sur l'ensemble des fiches de paie. Le tableau~\ref{tab:types-anomalies} en donne les types. Les neuf types déclarés dans l'application ne sont pas tous produits par le registre : les trois derniers sont prévus par le catalogue de recommandations pour d'autres contrôles.

\begin{table}[H]
\centering
\small
\renewcommand{\arraystretch}{1.2}
\tablestyle
\begin{tabularx}{\textwidth}{|>{\raggedright\arraybackslash}p{4.8cm}|Y|}
\hline
\rowcolor{appMarine}\entete{Type d'anomalie} & \entete{Origine et signification} \\
\hline
\texttt{IBAN\_INVALIDE} & Règle métier : IBAN absent, mal formé ou à clé incorrecte \\
\hline
\texttt{BIC\_INVALIDE} & Règle métier : BIC renseigné mais mal formé \\
\hline
\texttt{MONTANT\_INVALIDE} & Règle métier : montant nul ou négatif \\
\hline
\texttt{DOUBLON\_PAIEMENT} & Règle métier : plusieurs fiches avec le même IBAN et le même montant \\
\hline
\texttt{IBAN\_PARTAGE} & Règle métier : un même IBAN utilisé par plusieurs salariés \\
\hline
\texttt{MONTANT\_SUSPECT} & Module d'IA : montant atypique, avec ses raisons \\
\hline
\texttt{NET\_SUPERIEUR\_BRUT},\newline \texttt{ERREUR\_XSD},\newline \texttt{CTRL\_SUM\_INCOHERENT} & Types prévus par le catalogue de recommandations, non alimentés par la détection actuelle \\
\hline
\end{tabularx}
\caption{Types d'anomalies du registre}
\label{tab:types-anomalies}
\end{table}

Une anomalie a un statut de suivi : \textbf{ouverte}, \textbf{corrigée} ou \textbf{ignorée}. Les deux profils métier peuvent la marquer corrigée ou l'ignorer ; l'ignorer exige un commentaire, qui garde la trace du motif. La détection peut être relancée à la demande. Elle réconcilie le registre avec l'état actuel des fiches : une nouvelle anomalie est créée, une anomalie qui n'est plus détectée passe à \og corrigée\fg{}, et une anomalie corrigée à la main qui réapparaît est rouverte. Les anomalies ignorées ne sont pas modifiées. Si l'IA est indisponible lors d'une détection, les anomalies d'origine IA restent ouvertes au lieu d'être classées à tort comme corrigées.

Chaque anomalie reçoit à sa création une \textbf{recommandation} tirée d'un catalogue par type. Le texte s'appuie sur les valeurs réelles : matricule, IBAN masqué, montant et période. Par exemple, pour un IBAN invalide, la recommandation conseille de demander un RIB à jour au salarié, de corriger l'IBAN puis de relancer la génération. En cas d'incident du moteur de recommandation, un message générique est utilisé afin de ne jamais interrompre la détection. La figure~\ref{fig:ecran-anomalies} montre le registre.

Le registre est partagé par les deux profils métier, mais l'Expert RH ne voit que les anomalies des fiches de ses sociétés et ne peut modifier le statut que de celles-ci ; le Manager voit l'ensemble.

\subsection{Soumission et décision}

La \textbf{soumission}, réservée à l'Expert RH, exige quatre conditions : l'accès à la société du fichier, la conformité au schéma, l'absence d'erreur de contrôle métier et le statut \og Brouillon\fg{}. Elle enregistre l'auteur et la date, trace l'événement dans le journal d'audit et fait passer le fichier au statut \og Soumis à validation\fg{}. Lorsque le fichier comporte une erreur métier, le serveur refuse la soumission et l'interface masque le bouton en indiquant la raison (figure~\ref{fig:ecran-synthese}).

Le Manager retrouve les fichiers soumis dans une liste dédiée, triée du plus ancien au plus récent. Il ouvre le résultat complet, puis valide ou rejette (figure~\ref{fig:ecran-decision}). Le rejet exige un motif, qui est conservé et affiché ensuite à l'Expert RH. Dans les deux cas, le statut, l'auteur de la décision et sa date sont enregistrés et le journal d'audit en garde la trace. Seul un fichier soumis peut être validé ou rejeté.

Le statut conditionne aussi le téléchargement. Un fichier validé est téléchargé sous son nom d'origine. Tout autre téléchargement reçoit le préfixe \og NON\_VALIDE\fg{}, afin qu'un fichier de contrôle ne soit pas confondu avec un fichier destiné à la banque. SEPAPaie ne transmet pas le fichier à la banque : cette étape reste hors du périmètre.

\subsection{Historique et supervision}

L'historique recense les fichiers avec leur date, leur société, leur période, le nombre de salariés, le total, la conformité XSD, le résultat de l'IA, le statut et l'auteur. Il se filtre par statut et par recherche libre. Un Expert RH n'y voit que les fichiers de ses sociétés ; le Manager les voit tous. L'administrateur dispose d'une \textbf{vue de supervision} en lecture seule : il suit l'activité, mais ne voit ni les matricules, ni les motifs de rejet, ni les acteurs du circuit, et ne dispose d'aucune action (figure~\ref{fig:ecran-supervision}).

\subsection{Interfaces réalisées}

La figure~\ref{fig:ecran-preparation} présente la préparation d'un virement par un Expert RH. La liste affiche les contrôles de chaque paiement : la fiche EMP002 du jeu de démonstration porte un IBAN dont la clé est incohérente, ce qui est signalé en rouge.

\begin{figure}[H]
\centering
\includegraphics[width=0.95\textwidth]{fig6_preparation}
\caption{Préparation d'un virement de masse (profil Expert RH)}
\label{fig:ecran-preparation}
\end{figure}

La figure~\ref{fig:ecran-synthese} montre la synthèse du résultat. Les cartes résument les quatre dimensions du contrôle : validation XSD, contrôles métier, contrôle IA et statut du fichier. Ici, l'erreur métier bloque la soumission, et le message l'explique. La figure~\ref{fig:ecran-controles} détaille cette erreur dans l'onglet dédié.

\begin{figure}[H]
\centering
\includegraphics[width=0.95\textwidth]{fig6_synthese}
\caption{Synthèse du résultat : soumission bloquée par une erreur métier}
\label{fig:ecran-synthese}
\end{figure}

\begin{figure}[H]
\centering
\includegraphics[width=0.95\textwidth]{fig6_controles_metier}
\caption{Onglet des contrôles métier}
\label{fig:ecran-controles}
\end{figure}

Après la soumission, le Manager ouvre le résultat du fichier : les boutons de validation et de rejet apparaissent sous le résultat (figure~\ref{fig:ecran-decision}). La figure~\ref{fig:ecran-supervision} montre la vue de supervision de l'administrateur, sans actions ni détail par salarié.

\begin{figure}[H]
\centering
\includegraphics[width=0.95\textwidth]{fig6_decision}
\caption{Fichier soumis : décision du Manager}
\label{fig:ecran-decision}
\end{figure}

\begin{figure}[H]
\centering
\includegraphics[width=0.95\textwidth]{fig6_supervision}
\caption{Vue de supervision de l'historique (profil administrateur)}
\label{fig:ecran-supervision}
\end{figure}

Enfin, la figure~\ref{fig:ecran-anomalies} présente le registre des anomalies avec leurs recommandations et leurs actions de suivi.

\begin{figure}[H]
\centering
\includegraphics[width=0.95\textwidth]{fig6_anomalies}
\caption{Registre des anomalies (profil Manager)}
\label{fig:ecran-anomalies}
\end{figure}

\section{Conclusion}

Le sprint 3 a assemblé les briques précédentes en un parcours de virement de masse complet : préparation guidée, contrôles métier, génération et validation, suivi des anomalies avec recommandations, puis circuit de décision fondé sur les quatre yeux. Les fichiers non conformes ou comportant des erreurs de données ne peuvent pas être soumis, et chaque étape est tracée. Le chapitre suivant présente le sprint 4, consacré à la détection d'anomalies par intelligence artificielle, à son intégration et à son évaluation.



\section{Bilan des tests automatisés}

Les trois sprints sont couverts par 132 tests unitaires du backend, tous réussis (tableau~\ref{tab:tests-s123}). Ils utilisent des substituts pour la messagerie et le service Python afin d'éprouver les cas d'échec sans dépendance externe. Ils portent notamment sur les règles qui conditionnent la sécurité : un mot de passe correct n'ouvre jamais de session sans code, un service facial indisponible ne laisse jamais passer, un fichier non conforme ou comportant une erreur métier ne peut ni être soumis ni être téléchargé, et la validation exige un autre utilisateur que l'auteur. Ces tests vérifient la logique du backend ; ils ne remplacent pas la recette de l'application complète présentée au chapitre~\ref{ch:sprint5}.

\begin{table}[H]
\centering
\footnotesize
\renewcommand{\arraystretch}{1.1}
\tablestyle
\begin{tabularx}{\textwidth}{|>{\raggedright\arraybackslash}p{3.9cm}|Y|>{\centering\arraybackslash}p{1.1cm}|>{\centering\arraybackslash}p{1.3cm}|}
\hline
\rowcolor{appMarine}\entete{Domaine} & \entete{Comportements vérifiés} & \entete{Sprint} & \entete{Tests} \\
\hline
Authentification et code à usage unique & Session ouverte seulement après un code correct ; mot de passe, code ou visage erroné refusé ; code expiré, trop de tentatives, renvoi trop rapide, échec d'envoi & 1 & 15 \\
\hline
Inscription et reconnaissance faciale & Compte créé inactif ; adresse déjà utilisée ; approbation et rejet ; empreinte enregistrée ; visage différent, compte inactif ou sans visage refusé & 1 & 9 \\
\hline
Sociétés, affectations et droits & IBAN invalide refusé et masqué ; réaffectation avec historique ; droits propres à chaque profil, anciens rôles refusés & 1 & 15 \\
\hline
Fiches de paie & Société obligatoire et active ; approbation refusée pour un IBAN, un BIC ou un montant invalide ; approbation tracée & 2 & 14 \\
\hline
Schéma, génération et validation XSD & Analyse du schéma officiel ; nombre de virements et somme de contrôle ; versions du message ; collecte de toutes les erreurs avec leur ligne ; XML mal formé & 2 & 19 \\
\hline
Parcours de contrôle et circuit de validation & Paiements approuvés uniquement ; fichier non conforme conservé mais non téléchargeable ; soumission, validation et rejet selon les quatre yeux ; motif de rejet obligatoire & 3 & 36 \\
\hline
Contrôles de données et anomalies & Messages métier par paiement ; doublon en avertissement ; clé modulo~97 ; commentaire pour ignorer une anomalie ; anomalies d'IA conservées si l'IA est indisponible & 3 & 20 \\
\hline
Accès par société et supervision & Un Expert RH ne voit que ses sociétés ; vue de supervision limitée aux informations agrégées & 3 & 4 \\
\hline
\rowcolor{appClair}\textbf{Total} & & & \textbf{132} \\
\hline
\end{tabularx}
\caption{Tests automatisés des sprints 1 à 3}
\label{tab:tests-s123}
\end{table}
